Inside Partial Pooling the indifference of the marginal hider pins down the belief $e^*$, and with it the number of silent informed sellers. Renovators therefore come out of the stock of sellers who would have disclosed a low rating. A subsidy buys renovation but not disclosure until that stock is exhausted: the premium does not move, and the share of the remaining informed low-quality sellers who hide rises. The stock is large. Disclosed poor ratings are \attGeDisclosedShareApt{} percent of informed low-quality apartment listings and \attGeDisclosedShareHouse{} percent of house listings, so up to that share can renovate at the partial-equilibrium subsidy $S^*$ without any effect on hiding. Beyond it the remaining sellers hide and the threshold falls, so the subsidy needed for full take-up exceeds $S^*$ by $\psi b(\bar\theta-e^*)$, which is at most the width of the Partial Pooling band (Corollary~\ref{cor:att_band}). At the estimates this correction averages \attGeCorrPctBApt{} percent of $b$ in apartments (\euro{}\attGeCorrEurApt{} per m\textsuperscript{2}) and \attGeCorrPctBHouse{} percent in houses (\euro{}\attGeCorrEurHouse{}), so the partial-equilibrium subsidy is the general-equilibrium one to within about one percent of the rating premium. Uniqueness follows because renovation decisions are strategic substitutes through the silent pool: each renovation lowers the number of disclosed ratings and raises the gain from hiding for those left, which lowers the threshold for the next renovator. The proposition has an implication that does not rely on any calibrated quantity: in markets where a subsidy reaches more informed low-quality sellers, the registry hiding rate among those who remain should rise while the non-disclosure premium stays flat. Testing it requires programme take-up by commune and registry data after the subsidy.
